\section{从 tiny MLP 走向 decoder-only 模型}

Reference decoder 和 GPT-2 reference path 是 LCQI 从教学闭环走向真实模型导入的关键层。这里的代码更长，是因为真实模型不只有矩阵乘法：它还需要 tokenizer、位置、归一化、QKV 拆分、attention mask、KV cache、激活函数、lm head、权重命名映射、dtype 转换和坏资产拒绝。

\section{Tiny Reference Decoder}

\filepath{reference\_decoder.cpp} 实现 embedding、RMSNorm、Q/K/V、RoPE、KV cache、attention、SwiGLU、final norm、logits 和 sampler，并在关键位置生成 trace checkpoint。读者应按 checkpoint 顺序读，而不是从最终 logits 倒推。

\lstinputlisting[language=C++,caption={src/reference\_decoder.cpp}]{../../../cpu-volume-3/labs/linux_cpu_inference/src/reference_decoder.cpp}

\section{Tokenizer 与 Safetensors 实现}

\filepath{tokenizer.cpp} 把 prompt 变成 token id，并固定 unknown token 行为。读者应检查 BOS/EOS 是否稳定插入，UNK 是否按合同处理。

\lstinputlisting[language=C++,caption={src/tokenizer.cpp}]{../../../cpu-volume-3/labs/linux_cpu_inference/src/tokenizer.cpp}

\filepath{safetensors.cpp} 解析 header、metadata、dtype、shape 和 data offsets。坏 dtype、坏 shape、坏 offset 和 payload 太短都应该被拒绝。

\lstinputlisting[language=C++,caption={src/safetensors.cpp}]{../../../cpu-volume-3/labs/linux_cpu_inference/src/safetensors.cpp}

\section{GPT-2 Reference Path}

\filepath{gpt2\_reference.cpp} 包括 JSON 配置读取、byte-level BPE、safetensors F32/F16/BF16 张量读取、HuggingFace 权重名前缀解析、LayerNorm、packed QKV、causal attention、GELU、tied lm head、full-prefix greedy generation 和 cached-KV generation。它是正确性优先的 reference path，暂不追求吞吐。

\lstinputlisting[language=C++,caption={src/gpt2\_reference.cpp}]{../../../cpu-volume-3/labs/linux_cpu_inference/src/gpt2_reference.cpp}

\filepath{gpt2\_cli.cpp} 负责把模型目录、prompt、\code{max\_new\_tokens} 和 benchmark 选项接到 GPT-2 reference path。无参数时运行 tiny GPT-2 fixture；传入 HuggingFace 模型目录时加载真实资产。

\lstinputlisting[language=C++,caption={src/gpt2\_cli.cpp}]{../../../cpu-volume-3/labs/linux_cpu_inference/src/gpt2_cli.cpp}
